\newif\ifglobal

%%%%%%%%%%%%%%%%%%%%%%%
%%% Compile to view %%%
%%%%%%%%%%%%%%%%%%%%%%%

%%%%%%%%%%%%%%%%%%%%%%%%%%%%%%%%%%%%%%%%%%%%%%%%%%%
%%% Readme:                                     %%%
%%% https://www.overleaf.com/read/bwdjvfjksvjy/ %%%
%%%%%%%%%%%%%%%%%%%%%%%%%%%%%%%%%%%%%%%%%%%%%%%%%%%

% >>> M?: marks a module setting number or important notice

% >>> M4: Set {./relative-path} to external files for this module <<<
\def \modulefiles {./25-RNG}
% To add an external file, use:
% (Replace '---' with actual filename)
% '\input{\modulefiles/tables/---.tex}' for tables
% '\includegraphics{\modulefiles/figures/---}' for figures
% '\subfile{\modulefiles/---}' for register files


% >>> M1: This module is in global project? <<<
% 'YES' - uncomment; 'NO' - comment out
\globaltrue

\ifglobal
    \documentclass[main\_\_CN.tex]{subfiles}

\else
    % Font "FandolSong-Regular" used by ctex does not have GSUB table.
    % It causes warnings that do not affect compilation and can be ignored.
    \PassOptionsToPackage{quiet}{fontspec}

    \documentclass[openany, 10pt]{book}

    % >>> M2: In ./00-shared/preamble-trm-module, also find
    %         settings for visibility of todo notes, watermark <<<
    \usepackage{preamble-trm-module}


    % Enable Chinese typesetting
    \usepackage{ctex}

    % >>> M3: Temporary labels in Chinese
    %         you can add your own labels there <<<
    \myexternaldocument{./00-shared/temp-labels-cn}

\fi %\ifglobal


\setcounter{secnumdepth}{4}


% >>> M5: If this module needs any updates in
% './00-shared/preamble-trm-module.sty'
% (include additional package, etc.), contact your module PM
% or create the file './00-shared/changelog.tex' make the
% necessary updates yourself and document them in this file <<<


\begin{document}


% Sets language into which pre-defined bits of text are translated
\selectlanguage{chinese}

% Do not delete this block
\ifglobal
% Add the following front matter if
% this module is outside of global project
\else
    \InputIfFileExists{readme}{}{}
    {\let\clearpage\relax \listoftodos}
    {\let\clearpage\relax \listofdonetodos}
    \tableofcontents
    %\listoftables
    %\listoffigures
\fi


%%%%%%%%%%%%%%%%%%%%%%%%%
%%% TRM Module Begins %%%
%%%%%%%%%%%%%%%%%%%%%%%%%

\hypertarget{rng}{}
\chapter{随机数发生器 (RNG)}
\label{mod:rng}


\section{概述}

\chipname{} 内置一个真随机数发生器，其生成的 32 位随机数可作为加密等操作的基础。

\section{主要特性}

\chipname{} 的随机数发生器可通过物理过程而非算法生成真随机数，所有生成的随机数在特定范围内出现的概率完全一样。

\section{功能描述}\label{sec:rng-func-descr}

系统可以从随机数发生器的寄存器 \hyperref[regdesc:LPPERIRNGSYNCDATAREG]{LPPERI\_RNG\_DATA\_SYNC\_REG} 中读取随机数，每个读到的 32 位随机数都是真随机数，噪声源为系统中的\textbf{热噪声}、\textbf{异步时钟}和\textbf{异步计数器}。

\begin{itemize}
    \item \textbf{热噪声}可以来自 SAR ADC 或高速 ADC 或两者兼有。当芯片的 SAR ADC 或高速 ADC 工作时，就会产生比特流，并通过异或 (XOR) 逻辑运算作为随机数种子进入随机数生成器。
    \item RC\_FAST\_CLK 是一种\textbf{异步时钟源}，会产生电路亚稳态。这种亚稳态也可以作为随机数种子\hyperref[others:rng-note-clock-seed]{$^2$}，进入随机数生成器。
    \item \textbf{异步计数器}是利用异步时钟源进行计数，异步时钟源会产生电路亚稳态，可以作为随机数种子进入随机数生成器。异步计数器分为 2 种：
    \begin{itemize}
        \item 第一种是直接利用片上的 \textbf{\nameref{mod:rtctimer} (RTC\_TIMER)} 模块产生的计数值，并通过异或 （XOR） 逻辑运算作为随机数种子进入随机数生成器；
        \item 第二种是使用 BUF CELL (缓存单元）实现的环形振荡器 (\textbf{BUF\_CHAIN}) 生成的计数值，通过异或 (XOR) 逻辑运算作为随机数种子进入随机数生成器。
    \end{itemize}
    这两种异步计数器可以各自独立工作，也可以多个一起工作。
\end{itemize}

\begin{figure}[H]
    \centering
    \includegraphics[width=0.9\textwidth]{\modulefiles/figures/25-RNG-noise-source}
    \caption{噪声源}
    \label{fig:rng-noise-source}
\end{figure}

当 SAR ADC 产生噪声时，随机数生成器在每个 SAR ADC 采样周期内获得 2 位的熵。SAR ADC 的采样周期通常为 66 kHz，由 SAR ADC 配置决定。因此，为了获得最大的熵值，建议读取 \hyperref[regdesc:LPPERIRNGSYNCDATAREG]{LPPERI\_RNG\_DATA\_SYNC\_REG} 寄存器时的速率不超过 66 kHz。

当高速 ADC 打开时，每个 APB 时钟周期（通常为 80 MHz）内，随机数发生器将获得 2 位的熵。因此，为了获得最大的熵值，建议读取 \hyperref[regdesc:LPPERIRNGSYNCDATAREG]{LPPERI\_RNG\_DATA\_SYNC\_REG} 寄存器时的速率不超过 5 MHz。

当 BUF\_CHAIN 打开时，每个 RC\_FAST\_CLK 可获取 1 位的熵。

我们在仅打开高速 ADC 的状态下，以 5 MHz 的速率从 \hyperref[regdesc:LPPERIRNGSYNCDATAREG]{LPPERI\_RNG\_DATA\_SYNC\_REG} 读取了 2 GB 的数据样本，并使用 Dieharder 随机数测试套件（版本 3.31.1）对样本进行了测试。最终，样本通过了所有测试。


\section{编程指南}

要开启随机数生成器，需向 \hyperref[fielddesc:LPPERIRNGCKEN]{LPPERI\_RNG\_CK\_EN} 写 1 开启随机数生成器时钟。使用 \nameref{mod:ecdsa} 时，无论 \hyperref[fielddesc:LPPERIRNGCKEN]{LPPERI\_RNG\_CK\_EN} 配置如何，随机数生成器时钟都会自动启动。

在使用 \chipname{} 的随机数生成器时，如果熵源没有开启，可能会导致产生伪随机数，应注意避免。其中，
\begin{itemize}
    \item SAR ADC 受控于 DIG ADC 控制器。详见章节 \ref{mod:adc} \textit{\nameref{mod:adc}}。
    \item 高速 ADC 在 Wi-Fi 或蓝牙开启时自动打开。
    \item RC\_FAST\_CLK 处于常开状态。
    \item 异步计数器
        \begin{itemize}
            \item RTC 定时器 (RTC\_TIMER) 通过置位 \hyperref[fielddesc:LPPERIRTCTIMEREN]{LPPERI\_RTC\_TIMER\_EN} 开启。
            \item 环形振荡器 (BUF\_CHAIN) 通过置位 \hyperref[fielddesc:LPPERIRNGSAMPLEENABLE]{LPPERI\_RNG\_SAMPLE\_ENABLE} 开启。
        \end{itemize}
\end{itemize}
\newpage
\begin{tiplisting}
\vspace{-2.5em}
\begin{enumerate}
    \label{others:rng-note-clock-seed}
    \item 在 Wi-Fi 开启时，极端情况下高速 ADC 有读值饱和的可能，这会降低熵值。因此，建议在 Wi-Fi 开启时，同时打开 SAR ADC 产生随机数。
    \item RC\_FAST\_CLK 时钟和异步计数器仅可以提高随机数发生器的墒值。然而，为了保证随机数发生器可以获得足够大的熵值，在使用随机数发生器时应至少保证 SAR ADC 或高速 ADC 处于工作状态。
\end{enumerate}
\end{tiplisting}

当需要获取多个随机数时，可连续读取\hyperref[regdesc:LPPERIRNGSYNCDATAREG]{LPPERI\_RNG\_DATA\_SYNC\_REG} 寄存器的值。但是在读取寄存器时，注意控制速率不要超过章节 \ref{sec:rng-func-descr} 的介绍熵的产生速率。如果配置多个熵源，建议不超过最低的噪声源产生熵的速率。

%%%%%%%%%%%%%%%%%%%%%%%%
%%% Register Summary %%%
%%%%%%%%%%%%%%%%%%%%%%%%

\newpage
\hypertarget{rng-reg-summ}{}
\section{寄存器列表}
\label{sec:rng-reg-summ}

本小节的所有地址均为相对于 随机数生成器 基地址的地址偏移量（相对地址），具体基地址请见章节 \ref{mod:sysmem} \textit{\nameref{mod:sysmem}} 中的表 \ref{tab:sysmem-base-address}。

请查看章节 \hyperref[glossary-access-types]{\textit{寄存器的访问类型}}，了解“访问”列缩写的含义。

\begin{longtable}{ | p{6cm} | p{7.2cm} | p{1.3cm} | p{1cm} | }
\hline\rowcolor{lightgray}
\textbf{名称} & \textbf{描述} &  \textbf{地址} & \textbf{访问} \\ \hline
\endhead
\hyperref[regdesc:LPPERIRNGCFGREG]{LPPERI\_RNG\_CFG\_REG} & RNG 配置寄存器 & 0x0024 & 可变 \\ \hline
\hyperref[regdesc:LPPERIRNGSYNCDATAREG]{LPPERI\_RNG\_DATA\_SYNC\_REG} & RNG 数据同步寄存器 & 0x0028 & RO \\ \hline
\end{longtable}



%%%%%%%%%%%%%%%%%
%%% Registers %%%
%%%%%%%%%%%%%%%%%


\newpage
\section{寄存器}

本小节的所有地址均为相对于 随机数生成器 基地址的地址偏移量（相对地址），具体基地址请见章节 \ref{mod:sysmem} \textit{\nameref{mod:sysmem}} 中的表 \ref{tab:sysmem-base-address}。

关于保留 (reserved) 域的处理，请查看章节 \hyperref[programming-reserved-field]{\textit{如何配置寄存器的保留域}}。


\begin{register}{H}{LPPERI\_RNG\_CFG\_REG}{0x0024}\label{regdesc:LPPERIRNGCFGREG}
\regfield{(reserved)}{20}{12}{00000000000000000000}%
\regfield{LPPERI\_RTC\_TIMER\_EN}{2}{10}{11}%
\regfield{(reserved)}{9}{1}{000000000}%
\regfield{LPPERI\_RNG\_SAMPLE\_ENABLE}{1}{0}{{0}}%
\reglabel{Reset}\regnewline%
\begin{regdesc}\begin{reglist}
\label{fielddesc:LPPERIRNGSAMPLEENABLE}\item [LPPERI\_RNG\_SAMPLE\_ENABLE] 配置是否使能 BUF\_CHAIN。\\0：关闭\\ 1：使能\\  (R/W)
\label{fielddesc:LPPERIRTCTIMEREN}\item [LPPERI\_RTC\_TIMER\_EN] 配置是否在 CRC 前使能 RTC 定时器。\\
0：关闭\\
1：使能\\ (R/W)
\end{reglist}\end{regdesc}
\end{register}


\begin{register}{H}{LPPERI\_RNG\_DATA\_SYNC\_REG}{0x0028}\label{regdesc:LPPERIRNGSYNCDATAREG}
\regfield{LPPERI\_RNG\_SYNC\_DATA}{32}{0}{00000000000000000000000000000000}%
\reglabel{Reset}\regnewline%
\begin{regdesc}\begin{reglist}
\label{fielddesc:LPPERIRNGSYNCDATA}\item [LPPERI\_RNG\_SYNC\_DATA] 表示 RNG 同步数据。 (RO)
\end{reglist}\end{regdesc}
\end{register}

\end{document}
